% Options for packages loaded elsewhere
\PassOptionsToPackage{unicode}{hyperref}
\PassOptionsToPackage{hyphens}{url}
%
\documentclass[
  10pt,
]{article}
\usepackage{amsmath,amssymb}
\usepackage{iftex}
\ifPDFTeX
  \usepackage[T1]{fontenc}
  \usepackage[utf8]{inputenc}
  \usepackage{textcomp} % provide euro and other symbols
\else % if luatex or xetex
  \usepackage{unicode-math} % this also loads fontspec
  \defaultfontfeatures{Scale=MatchLowercase}
  \defaultfontfeatures[\rmfamily]{Ligatures=TeX,Scale=1}
\fi
\usepackage{lmodern}
\ifPDFTeX\else
  % xetex/luatex font selection
\fi
% Use upquote if available, for straight quotes in verbatim environments
\IfFileExists{upquote.sty}{\usepackage{upquote}}{}
\IfFileExists{microtype.sty}{% use microtype if available
  \usepackage[]{microtype}
  \UseMicrotypeSet[protrusion]{basicmath} % disable protrusion for tt fonts
}{}
\makeatletter
\@ifundefined{KOMAClassName}{% if non-KOMA class
  \IfFileExists{parskip.sty}{%
    \usepackage{parskip}
  }{% else
    \setlength{\parindent}{0pt}
    \setlength{\parskip}{6pt plus 2pt minus 1pt}}
}{% if KOMA class
  \KOMAoptions{parskip=half}}
\makeatother
\usepackage{xcolor}
\usepackage[margin=0.9in]{geometry}
\usepackage{color}
\usepackage{fancyvrb}
\newcommand{\VerbBar}{|}
\newcommand{\VERB}{\Verb[commandchars=\\\{\}]}
\DefineVerbatimEnvironment{Highlighting}{Verbatim}{commandchars=\\\{\}}
% Add ',fontsize=\small' for more characters per line
\newenvironment{Shaded}{}{}
\newcommand{\AlertTok}[1]{\textcolor[rgb]{1.00,0.00,0.00}{\textbf{#1}}}
\newcommand{\AnnotationTok}[1]{\textcolor[rgb]{0.38,0.63,0.69}{\textbf{\textit{#1}}}}
\newcommand{\AttributeTok}[1]{\textcolor[rgb]{0.49,0.56,0.16}{#1}}
\newcommand{\BaseNTok}[1]{\textcolor[rgb]{0.25,0.63,0.44}{#1}}
\newcommand{\BuiltInTok}[1]{\textcolor[rgb]{0.00,0.50,0.00}{#1}}
\newcommand{\CharTok}[1]{\textcolor[rgb]{0.25,0.44,0.63}{#1}}
\newcommand{\CommentTok}[1]{\textcolor[rgb]{0.38,0.63,0.69}{\textit{#1}}}
\newcommand{\CommentVarTok}[1]{\textcolor[rgb]{0.38,0.63,0.69}{\textbf{\textit{#1}}}}
\newcommand{\ConstantTok}[1]{\textcolor[rgb]{0.53,0.00,0.00}{#1}}
\newcommand{\ControlFlowTok}[1]{\textcolor[rgb]{0.00,0.44,0.13}{\textbf{#1}}}
\newcommand{\DataTypeTok}[1]{\textcolor[rgb]{0.56,0.13,0.00}{#1}}
\newcommand{\DecValTok}[1]{\textcolor[rgb]{0.25,0.63,0.44}{#1}}
\newcommand{\DocumentationTok}[1]{\textcolor[rgb]{0.73,0.13,0.13}{\textit{#1}}}
\newcommand{\ErrorTok}[1]{\textcolor[rgb]{1.00,0.00,0.00}{\textbf{#1}}}
\newcommand{\ExtensionTok}[1]{#1}
\newcommand{\FloatTok}[1]{\textcolor[rgb]{0.25,0.63,0.44}{#1}}
\newcommand{\FunctionTok}[1]{\textcolor[rgb]{0.02,0.16,0.49}{#1}}
\newcommand{\ImportTok}[1]{\textcolor[rgb]{0.00,0.50,0.00}{\textbf{#1}}}
\newcommand{\InformationTok}[1]{\textcolor[rgb]{0.38,0.63,0.69}{\textbf{\textit{#1}}}}
\newcommand{\KeywordTok}[1]{\textcolor[rgb]{0.00,0.44,0.13}{\textbf{#1}}}
\newcommand{\NormalTok}[1]{#1}
\newcommand{\OperatorTok}[1]{\textcolor[rgb]{0.40,0.40,0.40}{#1}}
\newcommand{\OtherTok}[1]{\textcolor[rgb]{0.00,0.44,0.13}{#1}}
\newcommand{\PreprocessorTok}[1]{\textcolor[rgb]{0.74,0.48,0.00}{#1}}
\newcommand{\RegionMarkerTok}[1]{#1}
\newcommand{\SpecialCharTok}[1]{\textcolor[rgb]{0.25,0.44,0.63}{#1}}
\newcommand{\SpecialStringTok}[1]{\textcolor[rgb]{0.73,0.40,0.53}{#1}}
\newcommand{\StringTok}[1]{\textcolor[rgb]{0.25,0.44,0.63}{#1}}
\newcommand{\VariableTok}[1]{\textcolor[rgb]{0.10,0.09,0.49}{#1}}
\newcommand{\VerbatimStringTok}[1]{\textcolor[rgb]{0.25,0.44,0.63}{#1}}
\newcommand{\WarningTok}[1]{\textcolor[rgb]{0.38,0.63,0.69}{\textbf{\textit{#1}}}}
\setlength{\emergencystretch}{3em} % prevent overfull lines
\providecommand{\tightlist}{%
  \setlength{\itemsep}{0pt}\setlength{\parskip}{0pt}}
\setcounter{secnumdepth}{-\maxdimen} % remove section numbering
\usepackage{amsmath,amssymb,amsthm}
\ifLuaTeX
  \usepackage{selnolig}  % disable illegal ligatures
\fi
\usepackage{bookmark}
\IfFileExists{xurl.sty}{\usepackage{xurl}}{} % add URL line breaks if available
\urlstyle{same}
\hypersetup{
  pdftitle={The fixed-Reeb Monge--Ampere image: genuine approximation and the remaining realization gate},
  pdfauthor={Public-source reconstruction; unverified partial mathematical draft},
  hidelinks,
  pdfcreator={LaTeX via pandoc}}

\title{The fixed-Reeb Monge--Ampere image: genuine approximation and the
remaining realization gate}
\author{Public-source reconstruction; unverified partial mathematical
draft}
\date{5 October 2026}

\begin{document}
\maketitle

\subsection{Result and exact status}\label{result-and-exact-status}

The complete original P13.2 is \textbf{not solved by this document}. We
prove the following partial theorem on its entire domain, including
irrational Reeb fields. The theorem gives an intrinsic, finite-test
description of the \emph{finite-energy target class}, a necessary
condition for the genuine image, and quantitative genuine approximate
solutions. It does not identify that target class with the image.

Let \(k\) be any algebraically closed trivially valued
characteristic-zero field and let \((X,T;\xi)\) be a normal good
polarized affine cone of dimension \(n\geq1\). Put \(p=n-1\),
\(K=X_0^T\), and \(o\) for the vertex. Potentials below are restrictions
of the actual cone classes. Their reference \(\rho_\xi\) restricts to
zero on \(K\). The reference point \(v_{\mathrm{triv}}\) belongs to
\(K\). Define, for a Radon probability \(\mu\) on this fixed carrier, \[
 L_\xi(\mu)=\sup_{w\in H(\xi)}\left(E_\xi(w)-\int_Kw\,d\mu\right).                     \tag{1}
\] Every test in (1) is an explicit finite maximum of homogeneous
logarithmic branches with a primary list. In particular, the condition
\(L_\xi(\mu)<\infty\) involves neither an inverse potential nor unknown
finite-energy outputs.

\textbf{Partial theorem.} For every cone and every Reeb field in the
stated domain:

\begin{enumerate}
\def\labelenumi{\arabic{enumi}.}
\tightlist
\item
  The condition \(L_\xi(\mu)<\infty\) is equivalent to
  \(\mathcal E^1(\xi)\subset L^1(\mu)\). It is exactly the existing
  finite-energy measure condition \(E_\xi^\vee(\mu)<\infty\). The
  finite-energy measure class on \(K\) is independent of the Reeb field.
  Every genuine output belongs to this class.
\item
  Every such target has an \(H(\xi)\) maximizing sequence \(u_j\). Its
  \emph{genuine} measures \(\mathrm{MA}_\xi(u_j)\) converge weakly to
  \(\mu\). For \(f\) in the accepted fixed dense direction lattice
  \(D\), a potential with dual deficit \(d\) satisfies \[
   \left|\int_K f\,d(\mathrm{MA}_\xi(u)-\mu)\right|
   \leq d/s+C_f s\qquad(0<s\leq\tau),                                  \tag{2}
  \] where \(C_f=2pA(1+A\tau)^{p-1}\|f\|_K\) and \(A\) is an actual
  positive auxiliary slope. If \(d=0\), both signs give the equation on
  a genuine \(u\in\mathcal E^1\).
\item
  When \(p\geq1\), set \(a_p=p/4^{p+1}\). For any two finite-energy
  potentials and their deficits, \[
   I_\xi(u,v)\leq\frac{d_u+d_v}{2a_p}.                                  \tag{3}
  \] Thus every maximizing sequence is Cauchy modulo constants for the
  existing \(I_\xi\). After normalization at \(v_{\mathrm{triv}}\), its
  energies have a uniform finite lower bound. For every \(f\in D\) with
  actual auxiliary slope \(A\), \[
   \left|\int_Kf\,d(\mathrm{MA}_\xi(u)-\mathrm{MA}_\xi(v))\right|
   \leq\sqrt{2Ap\|f\|_K I_\xi(u,v)}.                                   \tag{4}
  \] This proves genuine weak continuity if an \(I_\xi\) limit exists in
  the actual \(\mathcal E^1/\mathbb R\).
\item
  The carrier upper-semicontinuous supremum property transfers through
  the accepted intrinsic root maps: it holds at one Reeb field if and
  only if it holds at any other. We prove the upper-semicontinuous
  transport, rather than infer it from order alone.
\item
  For \(n=1\), \(K\) is a point and the exact image is its sole
  probability. This special case is complete; it does not cover the
  original all-dimensional problem.
\end{enumerate}

The unresolved universal gate is realization of the Cauchy sequences in
(3) by actual finite-energy potentials. Neither completeness of that
actual space nor the singular normal rational envelope base needed by an
envelope route is established here. Formula (2) proves Euler--Lagrange
at a genuine maximizer; it does not manufacture that maximizer. Formula
(4) would make completeness an independent realization substitute,
without requiring an unbounded moving-Reeb measure limit. No
variational-attainment condition or output-closure condition is
presented as a solution of P13.2.

\subsection{1. Accepted input and
normalization}\label{accepted-input-and-normalization}

We use the explicitly accepted \emph{full mathematical P13.1 proof} as a
supporting theorem, rather than use operational claims in its canonical
notes. Its exact binding is
\nolinkurl{/workspace/research/deliverables/original-P13.1/P13.1-reviewed-argument.md},
SHA256
\nolinkurl{de75b9ae55fbeb63c605505a328ed306c321a8e674331827cde51f5cfad76609}.
The independently accepted review is a supplied binding, not a new
review of this document. No new envelope conclusion is attributed to
that theorem.

The parts used are its actual-class convexity, decreasing-net energy and
weighted mixed-MA continuity, positive mass-one mixed products,
fixed-ray integration by parts and Hodge index, intrinsic order roots,
finite-energy root comparison, rational quotient calibration, dense
direction lattice, and canonical outer differentiation. In particular,
\[
 \nu_\xi(u_1,\ldots,u_p)
   =V_\xi^{-1}dd^cu_1\cdots dd^cu_p\,dd^c\rho_\xi^+,
 \quad \rho_\xi^+=\max(\rho_\xi,0),                                    \tag{5}
\] and \[
 E_\xi(u)=\frac1{nV_\xi}\sum_{j=0}^p
       \int_Ku\,(dd^cu)^j(dd^c\rho_\xi)^{p-j}dd^c\rho_\xi^+.             \tag{6}
\] For finite-energy potentials, (6) denotes the extended energy,
obtained by genuine fixed-field decreasing limits; the mixed weighted
integrals agree with that limit. There are \(p\) horizontal factors and
one radial factor in (5). Its total mass is one. Constants have energy
equal to the constant on \(K\).

For a rational field \(\eta=s\eta_e\), the effective supported grading
has gcd one. With the actual degree-\(d\) Veronese finite map and
\(L_d\) on the normal projective quotient \(V=\operatorname{Proj}R\),
the accepted calibration is \[
 \dim V=p,
 \quad V_\eta=s^{-n}d^{1-n}(L_d^p),
 \quad\omega_\eta=L_d/(sd),\quad (\omega_\eta^p)=sV_\eta,
 \quad E_\eta=E_{L_d}/(sd).                                           \tag{7}
\] The one-unit radial Gauss mass, the degree-\(d\) pushforward, and the
\(n\) factors \(1/(sd)\) are retained. The relevant Hilbert coefficients
are \((L_d^p)m^p/p!\) in a single projective degree and
\(V_\eta m^n/n!\) in a full cone cutoff. The signed full-ring mean is
\(S\) with \(E=(n+1)S/n\), not \(E=S\). These conventions are not
replaced by projective raw mass or transverse volume.

We also need the reference identity \[
                    \mathrm{MA}_\xi(\rho_\xi)=\delta_{v_{\mathrm{triv}}}.                    \tag{8}
\] Here is its deduction from the accepted bridge. Rationally the
canonical reference is the zero-weight product model on the normal
integral quotient. Its central fiber has one component of multiplicity
one; the mixed component intersection is \((L_d^p)\). Its probability
measure is therefore the trivial-valuation Dirac mass. The radial Gauss
lift and the carrier homeomorphism in the accepted proof give exactly
\(v_{\mathrm{triv}}\) on the cone. For an irrational \(\xi\), rational
references converge locally uniformly off \(o\) to the actual
\(\rho_\xi\). Their genuine probabilities converge by the accepted
\textbf{bounded} current-continuity result. All the preceding reference
measures are the same Dirac mass on the fixed \(K\), so their limit is
(8). No unbounded potential moves its Reeb field in this argument.

On \(K\), \(u(v_{\mathrm{triv}})=\sup_Ku\) is finite for each actual
nontrivial PSH potential. Write \(Nu=u-u(v_{\mathrm{triv}})\). For
normalized potentials, \(u\leq0\) and \[
              J_\xi(u)=-E_\xi(u)\geq0.                               \tag{9}
\] Convex combinations on \(K\) always mean the actual off-carrier
combination with the reference included so that the equivariant slope
remains one.

\subsection{2. The finite-test dual energy and target
integrability}\label{the-finite-test-dual-energy-and-target-integrability}

First prove the elementary integral fact in the nonmetrizable setting.
If \(h_i\) is a decreasing bounded-above usc net on compact Hausdorff
\(K\) and \(\mu\) is Radon, then \[
                  \int h_i\,d\mu\downarrow\int h\,d\mu,
             \qquad h=\inf_i h_i.                                    \tag{10}
\] An initial tail is bounded above by one constant. For a bounded-above
usc function, Radon regularity expresses its extended integral as the
infimum of integrals of continuous majorants. This follows first for
bounded functions by approximating their level sets by open sets, and
then by truncation below. If \(g\) is such a majorant of \(h\), the open
sets \(\{h_i<g+\epsilon\}\) increase and cover \(K\). Compactness and
directedness give \(h_i\leq g+\epsilon\) eventually. This proves the
upper bound in (10); the reverse follows from \(h_i\geq h\). The same
argument covers \(h=-\infty\) on a positive-measure set. It is not a
sequence-only monotone-convergence assertion.

Uniform actual \(H\) approximation on \(K\) and energy continuity show
that the supremum in (1) equals the supremum over CPSH. If
\(L_\xi(\mu)<\infty\) and \(u_i\downarrow u\in\mathcal E^1\) is its
actual CPSH defining net, then \[
             \int_Ku_i\,d\mu\geq E_\xi(u_i)-L_\xi(\mu).
\] Decreasing energy continuity and (10) imply \[
          \int_Ku\,d\mu\geq E_\xi(u)-L_\xi(\mu)>-\infty.               \tag{11}
\] Since \(u\) is bounded above on \(K\), this is \(L^1(\mu)\)
integrability. Applying both continuity statements to the difference
gives \[
 L_\xi(\mu)=\sup_{u\in\mathcal E^1(\xi)}
                  \left(E_\xi(u)-\int_Ku\,d\mu\right)
             =E_\xi^\vee(\mu).                                      \tag{12}
\] If the supremum over CPSH is infinite, the supremum over
\(\mathcal E^1\) is infinite as well. Thus the equality of extended dual
energies has no hidden prior integrability assumption.

We prove the converse integrability criterion, including its needed
fixed-ray scaling. If \(u\leq0\) is normalized and finite-energy, put \[
 b_l=-\int_K u\,d\nu_\xi(u^{[l]},\rho_\xi^{[p-l]})\geq0,
               \qquad 0\leq l\leq p.
\] All terms are finite by accepted weighted mixed continuity. Equation
(8) gives \(b_0=0\), while (6) gives \(J(u)=n^{-1}\sum_lb_l\). The
restriction \(tu\) for \(0\leq t\leq1\) represents
\(u_t=tu+(1-t)\rho_\xi\) off \(K\). Expanding the actual positive mixed
product in (6), \[
 J(u_t)=\frac1n\sum_{l=1}^p
   \left[t\sum_{j=l}^p{j\choose l}t^l(1-t)^{j-l}\right]b_l
       \leq 2^n t^2J(u).                                             \tag{13}
\] For the bound, \(t^{l+1}\leq t^2\), \((1-t)^{j-l}\leq1\), and
\(\sum_{j=l}^p{j\choose l}={n\choose l+1}\leq2^n\). The identity first
holds for continuous potentials, and passes through actual fixed-field
decreasing nets to \(\mathcal E^1\). No nonlinear cross-ray linear axiom
is invoked. If \(p=0\), all normalized potentials have zero energy on
the one-point carrier.

Suppose now that every actual finite-energy potential is
\(\mu\)-integrable. There is a finite \(B\) such that \[
 \int_K(-u)\,d\mu\leq B
 \quad\text{for all normalized }u\in\mathcal E^1\text{ with }J(u)\leq1.                       \tag{14}
\] Otherwise choose \(u_j\) in that class with
\(\int(-u_j)\,d\mu\geq2^j\). The restriction
\(w_m=\sum_{j=1}^m2^{-j}u_j\) is an actual convex combination after
adding the remaining reference weight. It decreases, is normalized at
\(v_{\mathrm{triv}}\), and concavity gives \(E(w_m)\geq-1\). Its
decreasing actual PSH limit is nontrivial because its value at that
reference point is zero. Decreasing energy continuity gives a genuine
\(w\in\mathcal E^1\) with \(E(w)\geq-1\). But (10), or ordinary monotone
convergence for this sequence, gives \(\int(-w)\,d\mu=\infty\). This
contradicts the hypothesis, proving (14).

For \(J(u)\geq2^{-n}\) choose \(t=(2^nJ(u))^{-1/2}\leq1\). Equations
(13)--(14) give \(\int(-u)\,d\mu\leq B\sqrt{2^nJ(u)}\). For smaller
\(J(u)\), (14) applies directly. Consequently the uniform bound \[
 \int(-u)\,d\mu\leq B+B\sqrt{2^nJ(u)},\qquad
 L_\xi(\mu)\leq B+2^{n-2}B^2<\infty                                 \tag{15}
\] follows by maximizing \(-x+B\sqrt{2^nx}\) over \(x\geq0\). This
completes the claimed integrability equivalence, without applying the
printed projective Theorem 9.20 to cones.

Finite-energy targets form a convex weakly lower-semicontinuous dual
class; the finite dual-energy sublevels are weakly compact because (1)
is a supremum of continuous affine functions on the compact space of
Radon probabilities on \(K\). These are compactness statements about
measures, not about actual potentials.

Finally, the accepted normalized root bijection obeys
\(u/c\leq\mathcal R_{\xi\eta}u\leq u/C\leq0\) and is onto the actual
normalized \(\mathcal E^1\) class. Therefore integrability of every
source potential implies integrability of every target potential, and
the inverse gives the converse. Combined with (15), this proves
independence of the finite-energy measure class from \(\xi\) on the
\textbf{same carrier}. It makes no strong-topology assertion about
roots.

\subsection{3. Necessity and the actual tangent
inequality}\label{necessity-and-the-actual-tangent-inequality}

Fixed-ray concavity and its first variation imply, for continuous actual
potentials, \[
 E(w)\leq E(u)+\int_K(w-u)\,d\mathrm{MA}(u).                           \tag{16}
\] Indeed apply concavity to the actual segment \(u+t(w-u)\) and its
derivative at \(t=0\). Approximate both potentials by raw decreasing
CPSH nets. The accepted energy and \textbf{weighted} fixed-field
mixed-MA continuity pass (16) to all \(u,w\in\mathcal E^1\).

For \(\nu=\mathrm{MA}(u)\), all finite-energy weights are
\(\nu\)-integrable by the same accepted mixed calculus. Equation (16)
and (12) give the exact identity \[
              L_\xi(\nu)=E(u)-\int_Ku\,d\nu<\infty.                  \tag{17}
\] The inequality \(\leq\) follows by applying (16) to every competitor;
equality follows by using \(u\) in (12). This proves the necessary
finite-energy condition and variational optimality of every genuine
solution.

Wu A Proposition 7.3 LF1249 prints the rearranged inequality in the
opposite direction. The concave tangent used here is (16), proved
independently from the accepted calculus. We do not use that printed
line as a premise.

\subsection{4. Outer duality, both signs, and genuine approximate
equations}\label{outer-duality-both-signs-and-genuine-approximate-equations}

Use the accepted canonical functional \[
 F(g)=\sup_{w\in\mathrm{CPSH}(\xi),\ w|_K\leq g}E(w),\qquad
 F^\downarrow(h)=\inf_{g\in C(K),\ g\geq h}F(g).                        \tag{18}
\] It is finite on continuous obstacles, monotone, translation
equivariant and 1-Lipschitz. It agrees with actual PSH energy on that
actual domain; its outer extension is decreasing-net continuous on
bounded-above usc functions. We use its accepted differentiability,
without asserting envelope attainment.

For continuous \(g\), positivity of \(\mu\) and the definition of \(F\)
give \[
          F(g)-\int_Kg\,d\mu\leq L_\xi(\mu).                         \tag{19}
\] In fact, a feasible \(w\leq g\) satisfies
\(E(w)-\int g\,d\mu\leq E(w)-\int w\,d\mu\leq L\). Conversely CPSH
restrictions are continuous and satisfy \(F(w)=E(w)\), so \(L\) is also
the supremum of the left side of (19).

For any bounded-above usc \(h\) that is \(\mu\)-integrable, choose its
decreasing continuous-majorant net. Such a net exists on compact
Hausdorff \(K\): take all continuous majorants below a fixed upper
constant, ordered by pointwise reverse order, using minima to make it
directed. The pointwise infimum is \(h\), by upper semicontinuity and
the continuous separation of a point from a closed level set. Equations
(10), (18), and (19) yield \[
          F^\downarrow(h)-\int_Kh\,d\mu\leq L_\xi(\mu).               \tag{20}
\] No interchange of an uncontrolled supremum and an unbounded integral
has occurred.

Fix a genuine \(u\in\mathcal E^1\) and let \[
 d_u=L_\xi(\mu)-E(u)+\int_Ku\,d\mu\geq0.
\] The weight \(h=u+tf\) is \(\mu\)-integrable for every continuous
bounded \(f\). The accepted outer derivative, for \(f\in D\), gives \[
 \left|F^\downarrow(u+tf)-E(u)-t\int_Kf\,d\mathrm{MA}(u)\right|
   \leq C_f t^2\quad(|t|\leq\tau),                                  \tag{21}
\] with the explicit constant in (2). Combining (20) and the lower side
of (21), \[
 t\int_Kf\,d(\mathrm{MA}(u)-\mu)\leq d_u+C_f t^2.
\] Taking \(t=s\) and \(t=-s\) proves (2). This is both-sign
Euler--Lagrange on a genuine potential if \(d_u=0\): send
\(s\downarrow0\), then use density of \(D\) and mass one to conclude
\(\mathrm{MA}(u)=\mu\). Conversely (17) gives \(d_u=0\) for a solution.
This is a proved optimality identity for an already genuine potential,
not a proposed intrinsic existence condition.

For every finite \(L\) choose \(u_j\in H(\xi)\) with
\(0\leq d_{u_j}<j^{-2}\), and normalize by a constant if desired.
Constants do not change deficits or MA. With \(\tau=1\) and \(s=1/j\),
(2) yields \[
 \left|\int_Kf\,d(\mathrm{MA}(u_j)-\mu)\right|
                         \leq(1+C_f)/j\qquad(f\in D).                 \tag{22}
\] For \(f\in C(K)\) choose \(f_m\in D\) uniformly close. The additional
error is at most \(2\|f-f_m\|_K\), because both measures are
probabilities. Choose \(m\) before \(j\). This proves weak convergence
on every continuous test, on any compact Hausdorff carrier and over
every field in the original domain. The maximizing sequence exists by
the real supremum in (1); no countable enumeration of the ground field
or all FS presentations is claimed.

More generally, the same argument applies to any maximizing net in
actual \(\mathcal E^1\). Its target integration uses (11) and (20). It
does not assert integration continuity along arbitrary nonmonotone
potential nets.

\subsection{5. The Cauchy property: an actual fixed-ray
calculation}\label{the-cauchy-property-an-actual-fixed-ray-calculation}

Assume \(p\geq1\). For continuous actual \(u,v\), let \(z=u-v\) and
define the nonnegative fixed-ray Dirichlet coefficients \[
 B_j=-V_\xi^{-1}\int z\,dd^cz\,(dd^cu)^j(dd^cv)^{p-1-j}dd^c\rho_\xi^+,
                      \qquad 0\leq j\leq p-1.                        \tag{23}
\] Nonnegativity is the accepted fixed-ray Hodge index. Integration by
parts and the difference-of-powers identity give \[
 I_\xi(u,v)=\int_K(u-v)\,d(\mathrm{MA}(v)-\mathrm{MA}(u))
                          =\sum_{j=0}^{p-1}B_j.                      \tag{24}
\] Writing \(g(t)=E(v+tz)\), the genuine derivative and the actual
product expansion give \[
 -g''(t)=p\sum_{j=0}^{p-1}{p-1\choose j}t^j(1-t)^{p-1-j}B_j.           \tag{25}
\] These are mixed products on the one fixed ray. The roots are absent
from this calculation.

The Green kernel on \([0,1]\), evaluated at \(1/2\), is
\(G(1/2,t)=\tfrac12\min(t,1-t)\). Consequently \[
 E((u+v)/2)-(E(u)+E(v))/2
       =\int_0^1G(1/2,t)(-g''(t))\,dt\geq a_p I_\xi(u,v).             \tag{26}
\] To check the coefficient, on \([1/4,3/4]\) the Green kernel is at
least \(1/8\), every \(t^j(1-t)^{p-1-j}\) is at least \(4^{-(p-1)}\),
each binomial is at least one, and the interval length is \(1/2\). Their
product with \(p\) is \(p/4^{p+1}=a_p\).

All expressions (23)--(24) can also be written as sums of accepted
weighted genuine mixed measures. Simultaneous raw decreasing CPSH
approximation therefore passes them, their nonnegativity, and (26) to
actual finite-energy potentials. The midpoint approximants decrease to
the actual midpoint. Thus no unbounded differential current is newly
postulated by the notation in (23).

The dual deficit of the actual midpoint equals \[
 d_{(u+v)/2}=(d_u+d_v)/2-
               \big(E((u+v)/2)-(E(u)+E(v))/2\big)\geq0.
\] Combining with (26) proves (3). This proves Cauchyness of every
maximizing sequence directly, without silently applying a projective or
synthetic completeness theorem.

Pair a normalized \(u\) with the reference. Its reference deficit is
\(d_\rho=L\). The two nonnegative tangent deficits from (16) sum to
\(I(u,\rho)\), and the reference one is \(-E(u)\) by (8). Hence \[
                   -E(u)\leq I(u,\rho)
                    \leq (d_u+L)/(2a_p).                             \tag{27}
\] The approximating normalized potentials thus have bounded finite
energy. This bound does not imply their compactness or their membership
after completion.

\subsection{6. A genuine MA continuity estimate if an actual limit
exists}\label{a-genuine-ma-continuity-estimate-if-an-actual-limit-exists}

For \(f\in D\), the accepted direction construction at this fixed
\(\xi\) supplies actual continuous slope-\(A\) cone psh functions
\(a,b\) with \(f=a-b\) on \(K\), where \(a/A,b/A\) belong to
CPSH\((\xi)\). Extend the difference with equivariant slope zero. This
is a common \textbf{linear fixed-ray} class, sufficient for the
following calculation; no synthetic theorem is applied.

For each horizontal product
\(M_j=(dd^cu)^j(dd^cv)^{p-1-j}dd^c\rho_\xi^+\), Hodge index and
Cauchy--Schwarz give \[
 \left|V_\xi^{-1}\int f\,dd^cz\,M_j\right|\leq\sqrt{Q_j(f)B_j},
 \qquad Q_j(f)=-V_\xi^{-1}\int f\,dd^cf\,M_j\geq0.
\] Since \(dd^cf=dd^ca-dd^cb\) and each positive product
\(V_\xi^{-1}dd^ca\,M_j\) or \(V_\xi^{-1}dd^cb\,M_j\) has mass \(A\),
positivity gives \[
                         Q_j(f)\leq2A\|f\|_K.                       \tag{28}
\] The mass assertion follows immediately by factoring out \(A\) and
applying the accepted mixed probability to \(a/A\) or \(b/A\) and the
remaining same-field slots. Summing the difference-of-powers identity
and applying the finite Cauchy--Schwarz inequality gives \[
 \left|\int f\,d(\mathrm{MA}(u)-\mathrm{MA}(v))\right|
 \leq\sqrt{2A\|f\|_K}\sum_j\sqrt{B_j}
 \leq\sqrt{2Ap\|f\|_K I(u,v)},
\] which is (4). For finite-energy \(u,v\), approximate decreasingly at
this \textbf{fixed} field. The accepted weighted continuity passes \(I\)
and the coefficients in (24); ordinary weak decreasing continuity passes
the bounded \(f\) tests. This proves the estimate on actual
\(\mathcal E^1\).

It follows that, if the Cauchy sequence in Section 5 has a limit \([u]\)
in the actual \(\mathcal E^1/\mathbb R\) for \(I\), its MA measures
converge weakly to \(\mathrm{MA}(u)\). By Section 4 they also converge
to \(\mu\), so \(\mathrm{MA}(u)=\mu\). The uniqueness assumed in the
task then applies. This explains precisely why \textbf{actual}
completeness would suffice as an independent realization route. It is a
reduction, not a theorem asserting completeness or a substitute image
condition.

This route does not need to pass target integrals through a nonmonotone
maximizing sequence: the actual equation follows from the two weak
measure limits. Once the genuine solution exists, its target integral is
finite by (11), and (17) and (21) give actual optimality and both-sign
Euler--Lagrange. The missing assertion is the existence of that actual
limit.

\subsection{7. Upper-semicontinuous envelope transport through
roots}\label{upper-semicontinuous-envelope-transport-through-roots}

Here is an analytic statement beyond the accepted order comparison, with
its proof. It concerns carrier envelopes and does not assume that an
envelope is psh.

Let \(\Omega\) be the torified analytic points off \(o\). With fixed
homogeneous generators \(h_i\) of weights \(\beta_i\), the accepted
action is jointly continuous and \[
 (v+t\eta)(\textstyle\sum f_\alpha)=\min_\alpha\{v(f_\alpha)+t\langle\eta,\alpha\rangle\}.
\] On \(\Omega\), every fixed-\(t\) map is a homeomorphism, with inverse
\(-t\). The reference and slice maps are \[
 \rho_\xi(v)=\max_i\frac{-v(h_i)}{\langle\xi,\beta_i\rangle},\qquad
 s_\xi(v)=v+\rho_\xi(v)\xi.
\] The map \(v\mapsto(s_\xi(v),\rho_\xi(v))\) is a homeomorphism
\(\Omega\simeq K\times\mathbb R\), with inverse \((x,r)\mapsto x-r\xi\).
Its continuity follows from the accepted action and the continuous
reference; the two identities verify the inverse. Thus for a function
\(h\) on \(K\) its equivariant extension
\(U(v)=h(s_\xi(v))+\rho_\xi(v)\) has usc regularization \[
                        U^*(v)=h^*(s_\xi(v))+\rho_\xi(v).              \tag{29}
\] This product-space equality is valid for extended \(-\infty\) values
and arbitrary nets. It does not require metrizability of \(K\).

Let \(\{u_i\}\) be any nonempty family of actual \(\xi\)-PSH potentials
with a uniform upper bound on their restrictions to \(K\). Set
\(U=\sup_i u_i\) on \(\Omega\). It is locally bounded above, is
\(\xi\)-equivariant, and retains the accepted orbit inequalities \[
 U(v)-Ct\leq U(v+t\eta)\leq U(v)-ct\quad(t\geq0),
 \] where \(c,C>0\) control all supported degree ratios. The same
inequalities hold for \(U^*\): take usc regularizations and use that
each fixed action map is a homeomorphism. Finiteness at one point of an
orbit forces finiteness everywhere on that orbit and a unique zero
parameter; an infinite orbit stays \(-\infty\). Define \(\mathcal R U\)
and \(\mathcal R U^*\) by those parameters, without claiming psh
membership.

The uniform lower slope proves the two identities \[
        \mathcal R U=\sup_i\mathcal R u_i,
        \qquad \mathcal R U^*=(\mathcal R U)^*.                        \tag{30}
\] For the first, put \(r=\sup_i\mathcal R u_i(v)\). Above \(r\), every
value is at most \(-c(t-r)\); below \(r\), some branch has a positive
value. If \(r=-\infty\), every original orbit is entirely infinite,
giving the asserted extended identity.

For the second, let \(r^*=\mathcal R U^*(v)\). If \(t>r^*\), then
\(U^*(v+t\eta)<0\). Upper semicontinuity gives a neighborhood where
\(U\) is strictly negative. Applying the fixed inverse action gives a
neighborhood of \(v\) on which \(\mathcal R U<t\), so
\((\mathcal R U)^*(v)\leq t\). If \(t<r^*\), the defining property of
usc regularization and \(U^*(v+t\eta)>0\) imply that every neighborhood
of \(v+t\eta\) meets \(\{U>0\}\). Apply the inverse action: every
neighborhood of \(v\) meets \(\{\mathcal R U>t\}\). Hence
\((\mathcal R U)^*(v)\geq t\). Taking arbitrary such \(t\) proves
equality. If \(r^*=-\infty\), the first argument works for every real
\(t\). There is no \(+\infty\) case, since the family is locally bounded
above.

Define the \textbf{carrier supremum property} at \(\xi\) to mean that
\((\sup_i u_i|_K)^*\) is the restriction of an actual \(\xi\)-PSH
potential for every such family. It is nontrivial at
\(v_{\mathrm{triv}}\). Root images of a uniformly upper-bounded family
are uniformly upper bounded: their values at \(v_{\mathrm{triv}}\) are
bounded by the same orbit slopes, and those values equal their carrier
suprema. If the property holds at \(\eta\), equations (29)--(30) make
the envelope of the root family an actual \(\eta\)-PSH potential. Apply
the accepted \textbf{raw} inverse root. It is an actual \(\xi\)-PSH
potential, and (29)--(30) identify its restriction with the source
carrier envelope. Reversing fields proves equivalence.

Only torified \(\Omega\) and \(K\) are involved in (29)--(30). An
envelope on the full link is not identified by an unjustified openness
assertion for torification. Actual potential membership in the last
paragraph follows from the accepted raw inverse and uniqueness of the
equivariant, torified extension.

At a rational field, the accepted normal quotient/torsor and residual
torification describe this carrier property as the invariant projective
envelope property. An independently proved unrestricted projective EP on
that actual singular normal quotient would suffice. It is not silently
supplied here. Universally this base cannot be reduced to smooth
quotients. To check this assertion, take a normal integral projective
variety \(V\) and an ample Cartier divisor \(D\). Its full section ring
is the intersection, inside \(k(V)(t)\), of \(k(V)[t]\) with the Gauss
valuation rings having weights \(v_F(f)+m\operatorname{ord}_F(D)\) on
\(ft^m\), for all prime divisors \(F\) of \(V\). Normality of \(V\)
identifies sections with these codimension-one inequalities. Every ring
in this intersection is integrally closed, so the section ring is
normal. Ampleness gives finite generation, and a high Veronese is
standard and has Proj equal to \(V\). Its affine section cone with the
scalar torus is therefore an original normal good cone whose rational
quotient is precisely \(V\) and whose residual torus is trivial. Thus
all normal singular projective bases occur in the original family.

No strong-topology continuity is inferred from (30). Although carrier
envelopes transfer in this precise sense, the accepted roots remain
nonlinear and do not by themselves transport \(I\) completeness.

\subsection{8. Why the printed singular source cannot close the gate as
supplied}\label{why-the-printed-singular-source-cannot-close-the-gate-as-supplied}

E is Lyuhui Wu, arXiv:2609.27610v1, the exact public primary cache with
PDF SHA256
\nolinkurl{11808a6c9ede9c0158b27e16700fc180a4ac8df12a4586a7178e1958e8f6bba8}.
Its extract SHA256 is
\nolinkurl{c8ffe91a631da67f6d3872a1a06ab8d15755aed2f926cdf326725ca3662cef7d}.
Coordinates below are inclusive one-based indices after UTF-8 byte
decoding and
\texttt{split(\textquotesingle{}\textbackslash{}n\textquotesingle{})}.

At LF766--768, Definition 3.15(2) calls \(v\) logarithmic-growth if
there is an FS function \(g\) with \(g\leq v+O(1)\). At LF769--774
regularizability requires a decreasing sequence of global FS functions
converging to \(v\). At LF1560--1562 Proposition 4.19 claims equality of
those classes on every smooth affine variety in equicharacteristic zero.
The proof at LF1569--1572 tries to get the opposite, upper-growth
estimate and invokes Lemma 3.25(1), whose hypothesis at LF952--957
already includes regularizability.

The lower-growth statement of Proposition 4.19 is actually false as
printed. This is an independently proved source diagnostic, not just an
observation about a missing citation. Take \(Y=\mathbb A^1_k\) and \[
             v(x)=\big(\max(\log|z|(x),0)\big)^2,
             \qquad g(x)=\max(\log|z|(x),0).                           \tag{31}
\] At \(z=0\) give both functions value zero. The function \(v\) is
continuous, finite everywhere on \(Y^{\mathrm{an}}\), and locally psh
under E Definition 3.8. To verify that exact definition, on any
neighborhood \(\{|z|<e^M\}\) with rational positive \(M\), let
\(t=-\log|z|\). Interpolate \(t^2\) on the successive dyadic
subdivisions of \([-M,0]\) by secants, and give the interpolants value
zero for \(t\geq0\), including \(t=+\infty\). These are continuous
convex rational PL functions on \((-M,+\infty]\), with slopes increasing
to zero. On the finite-coordinate part they are finite maxima of
rational affine branches and the zero branch; near \(+\infty\) they are
the constant zero. This finite local cover meets E's requirement that
the active affine expressions take finite values on their domains.
Refining the nested subdivision decreases the secant interpolants to
\((\max(-t,0))^2\). Composing with the tropicalization by \(z\) gives
precisely the decreasing local TConv--PL sequences in E Definitions
3.7--3.8, LF696--707. These neighborhoods cover \(Y^{\mathrm{an}}\).

The element \(z\) generates \(k[z]\), so \(g\) is an actual E-FS
function with rational data and constant branch zero. For every
\(r\geq0\), \(r\leq r^2+1/4\). Consequently \(g\leq v+1/4\) everywhere.
Thus \(v\) belongs to E's \textbf{printed lower-growth} class.

Nevertheless \(v\) cannot be a decreasing limit of global E-FS
functions. For \(t\geq0\), the Gauss seminorm of radius \(e^t\) on
\(k[z]\) is \(|\sum a_jz^j|_t=\max_j |a_j|e^{jt}\). Every nonzero
polynomial \(f\) has \(\log|f|_t=(\deg f)t\), since \(k\) is trivially
valued. Any E-FS function is a finite maximum of positive rational
multiples of these expressions plus rational constants and a rational
constant branch. It is bounded above by \(At+B\) for some finite \(A,B\)
on this Gauss ray. If a decreasing FS sequence converged to \(v\), its
first term would majorize \(v\), requiring \(t^2\leq At+B\) for all
\(t\geq0\). This is impossible. The counterexample is on a smooth affine
line and uses the exact rational/local definitions, without external
regularization results.

Replacing the lower-growth definition by an upper-growth definition
removes this particular counterexample. It does \textbf{not} prove the
corrected proposition or the singular envelope theorem. The global
extension step must then be justified without presupposing the very
regularizability it is trying to prove. The remaining trace, Runge,
strict-contact and singular-induction applications would still need
independent reconstruction with their antecedents, including
arbitrary-field descent. E Theorem 6.3 and 6.5 have
countable-dense-subfield hypotheses at LF1926--1929 and LF1938--1941,
while Theorem 7.1 removes that restriction through another descent
argument. These hypotheses cannot be suppressed.

As a separate normalization warning, the strict contact construction in
E Lemma 6.6 LF1972 prints \(\chi_k(t')>l(t)\), although contact at \(t\)
requires comparison with the affine function evaluated at \(t'\). We do
not import the printed contact lemma or claim that correcting the
displayed variable alone checks its construction. Similarly, the open
trace and singular induction bodies are retained for auditing, not
assumed valid. The conclusion here is limited: the supplied printed
chain is not a verified applicable singular EP proof. This document
neither disproves the final EP conclusion nor claims that no repair
exists.

\subsection{9. Dimension one and precise unresolved
conclusions}\label{dimension-one-and-precise-unresolved-conclusions}

The accepted normal effective-grading argument gives \(R=k[z]\) in cone
dimension one. Torification and the link condition force \(|z|=1\) on
\(K\), so \(K=\{v_{\mathrm{triv}}\}\). Every actual potential
restriction is a finite constant. The empty horizontal product and the
calibrated radial factor give
\(\mathrm{MA}(u)=\delta_{v_{\mathrm{triv}}}\), and
\(E(u)=u(v_{\mathrm{triv}})\). The only Radon probability on \(K\)
therefore has a genuine solution and dual energy zero, with uniqueness
modulo constants. There is no higher-dimensional consequence from this
case.

For \(n\geq2\), the strongest unconditional conclusion proved here is
(1)--(4), (12)--(17), (22), (27), and (30). It gives the following
concrete inventory.

\begin{itemize}
\tightlist
\item
  Finite-energy targets are determined by explicit \(H\) integrals and
  genuine energy, or equivalently by integrating the whole known
  \(\mathcal E^1\) domain. This is a description of the target class,
  not a proved description of the image.
\item
  The targets have genuine model approximate solutions, with a
  quantitative weak error and Cauchy actual potentials of bounded
  energy. Execution, finite arithmetic, and weak measure convergence do
  not create an actual inverse.
\item
  An actual \(I\) limit would solve the equation by (4). \textbf{Actual
  completeness on all normal cones is unproved here.} A formal
  completion cannot be substituted.
\item
  Carrier usc-sup stability transfers through roots, but \textbf{its
  singular normal rational base is unproved here}. If pursuing a
  compactness variational proof instead, the necessary
  target-integration continuity on the resulting nonmonotone potential
  limits would also have to be checked. The decreasing-net target
  continuity in Section 2 alone is insufficient for that different
  route.
\item
  The new singular source cannot be used by simply citing Theorem 7.1;
  Section 8 proves a counterexample to an indispensable proposition
  under its printed definition and leaves the corrected chain
  unverified.
\end{itemize}

The operator's given bi-Lipschitz embedding and uniqueness on its image
are fully consistent with this partial theorem. A Cauchy dense image can
remain a proper image until realization is established. No rational-only
or toric-only substitute covers the unchanged completion criteria.
Accordingly \texttt{candidate=false} and the status is partial, with an
explicit mathematical realization gap.

\subsection{10. Provenance, execution, and
rerun}\label{provenance-execution-and-rerun}

Only the input, the eleven supplied read-only primary caches and the
three explicitly accepted bindings were read. The accepted mathematical
body and the relevant literal primary bodies are retained in
\texttt{source-extracts/source-provenance.json}, with exact versions,
SHA256 checks, inclusive LF ranges, checked applications and concrete
limits. Byte-identical complete text copies
\texttt{source-extracts/A-full.txt} through \texttt{K-full.txt} retain
bodies beyond the selected excerpts. Retention is not a claim that all
their proofs were independently checked. This document applies no
synthetic theorem and no new singular EP theorem.

\texttt{input.json} is untouched; \texttt{input-original.json} is byte
identical. \texttt{assignment-original.txt} preserves the decoded
assignment, and \texttt{task-original.json} preserves the literal
task-object bytes. \texttt{preservation.json} records all four hashes.
\texttt{status.json} has exactly
\texttt{id,status,exact\_result,proof\_file,concrete\_gaps}.

The lightweight exact diagnostics in \texttt{checks.py} check the
coefficient bound in (13) on a finite rational grid, nested dyadic
secant inequalities used in (31), and the scalar factors in (7). They
are exhaustive only on their stated finite domains. They prove neither
universal mathematical correctness nor realization. No heavy
computation, worker, other model, network request, software
installation, account/runtime change, external message or publication is
used.

The existing Pandoc and pdfTeX tools generate \texttt{P13.2-paper.tex}
and compile twice. Separate \texttt{pandoc.stdout},
\texttt{pandoc.stderr}, \texttt{tex-pass-1.stdout},
\texttt{tex-pass-1.stderr}, \texttt{tex-pass-2.stdout},
\texttt{tex-pass-2.stderr}, \texttt{tex-pass-1.log}, and
\texttt{tex-pass-2.log} are retained with \texttt{environment.json} and
\texttt{build-report.json}. Successful compilation establishes
readability and execution only.

Rerun in this fresh workspace:

\begin{Shaded}
\begin{Highlighting}[]
\ExtensionTok{python}\NormalTok{ audit\_sources.py}
\ExtensionTok{python}\NormalTok{ checks.py}
\ExtensionTok{python}\NormalTok{ build.py}
\ExtensionTok{python}\NormalTok{ validate.py}
\end{Highlighting}
\end{Shaded}

The complete partial proof is this \texttt{P13.2-argument.md}, with
standalone \texttt{P13.2-paper.tex}. The final canonical notes preserve
the claims, proofs, diagnostic and exact gaps; none claims a complete
solution of the original P13.2.

\end{document}
